\documentclass[11pt]{article}
\usepackage[margin=1in]{geometry}
\usepackage{booktabs}
\usepackage{array}
\usepackage{multirow}
\usepackage{makecell}
\usepackage{amsmath}
\usepackage{amssymb}
\usepackage{graphicx}
\usepackage{caption}
\usepackage[numbers,sort&compress]{natbib}

\begin{document}

% Reference numbering is fixed by this list, not by where each work is first cited:
% 1-6 are the framing works and the metric definitions, 7 onwards follow the rows of
% Table~\ref{tab:scared_sota} from top to bottom. Remove this block to fall back to
% IEEE's default order-of-first-citation numbering.
\nocite{stoyanov2012surgical}
\nocite{allan2021scared}
\nocite{wang2025vggt}
\nocite{zhang2025monst3r}
\nocite{eigen2014depth}
\nocite{sturm2012benchmark}
\nocite{fang2020goodpractice}
\nocite{godard2019monodepth2}
\nocite{ozyoruk2021endosfm}
\nocite{shao2022afsfmlearner}
\nocite{zeinoddin2024dares}
\nocite{cui2024endodac}
\nocite{zeinoddin2025endofast3r}
\nocite{zhang2026endosfm3d}
\nocite{chang2026esrt}
\nocite{wang2024dust3r}


\section*{Training Losses}

We fine-tune the model with a multi-task loss that keeps VGGT's~\cite{wang2025vggt} original
supervision and adds
two terms specific to our method:
\begin{equation}
  \mathcal{L} = \lambda_{\text{cam}} \mathcal{L}_{\text{camera}}
              + \lambda_{\text{depth}} \mathcal{L}_{\text{depth}}
              + \lambda_{\text{smooth}} \mathcal{L}_{\text{smooth}},
  \qquad
  \lambda_{\text{cam}} = 5,\quad
  \lambda_{\text{depth}} = 1,\quad
  \lambda_{\text{smooth}} = 3 .
\end{equation}
Unlike the original formulation, the three terms are not on comparable scales here: the
camera term is up-weighted because pose is the target of this work, and the smoothness term
is a regulariser whose magnitude is a velocity rather than an error. The point-map and
tracking losses are not used, as the corresponding heads stay frozen.

\paragraph{Camera loss.}
The camera loss supervises the predicted cameras $\hat{g}_i$ against the ground truth $g_i$.
Each camera is encoded as $g_i = [\mathbf{t}_i, \mathbf{q}_i, f_i]$ --- translation,
rotation quaternion and (scalar) field of view --- and the three components are weighted
separately:
\begin{equation}
  \mathcal{L}_{\text{camera}}
  = \frac{1}{N} \sum_{i=1}^{N}
    \Big( w_{t}\big\|\hat{\mathbf{t}}_i - \mathbf{t}_i\big\|_{1}
        + w_{r}\big\|\hat{\mathbf{q}}_i - \mathbf{q}_i\big\|_{1}
        + w_{f}\big|\hat{f}_i - f_i\big| \Big),
  \qquad w_t = w_r = 1,\; w_f = 0.5 .
\end{equation}
Following the public VGGT implementation we use an $\ell_1$ loss here rather than the Huber
loss stated in the paper, as it was found to be more stable. The camera head refines its
prediction over several stages and every stage is supervised, with an exponentially decaying
weight $\gamma^{K-k}$, $\gamma = 0.6$, so that the final stage dominates. The translation term
is clamped at $100$, and frames whose ground-truth depth is empty are excluded from the sum,
since their pose is not observable.

\paragraph{Depth loss.}
The depth loss follows VGGT and implements the aleatoric-uncertainty loss: the per-pixel
discrepancy between the predicted depth $\hat{D}_i$ and the ground-truth depth $D_i$ is
weighted by a predicted \emph{confidence} map $C_i$, the depth head's second output channel,
which holds one strictly positive scalar per pixel and plays the role of an inverse
uncertainty. Let $\mathcal{V}_i$ be the pixels of frame $i$ with valid ground truth and
$\operatorname{mean}_{\mathcal{V}_i}$ the average over them; all operations below are
element-wise, so the pixel index can be dropped:
\begin{equation}
  \mathcal{L}_{\text{depth}}
  = \sum_{i=1}^{N} \operatorname*{mean}_{\mathcal{V}_i}
    \Big(
      C_i \big| \hat{D}_i - D_i \big|
      \;-\; \alpha \log C_i
      \;+\; \big\| \nabla \hat{D}_i - \nabla D_i \big\|_{1}
    \Big) ,
  \qquad \alpha = 0.2 ,
\end{equation}
where $|\cdot|$ is the absolute value of a scalar and $\nabla$ the spatial gradient in the image
plane. The first term makes the model confident where it is accurate, and $-\alpha \log C_i$
keeps the confidence from collapsing to zero; the two balance at
$C_i \propto 1 / \big| \hat{D}_i - D_i \big|$. The third term compares
the depth gradients rather than the depth values, which stops the model from buying a comparable
$\ell_1$ score with an over-smoothed prediction that erases tissue boundaries; it is evaluated at
four scales, obtained by subsampling the maps by a factor of $2^{s}$, $s = 0,\dots,3$, and
averaged over them, so that both fine edges and large-scale surface curvature are supervised.
Unlike the regression term it is deliberately not weighted by $C_i$, so the model cannot bypass
the structural constraint by lowering its confidence. Restricting the sums to $\mathcal{V}_i$
matters because endoscopic depth is sparse; the per-pixel terms are additionally trimmed at the
$0.98$ quantile to reject outliers.

\paragraph{Motion gate.}
The gate is not supervised by a loss of its own. A gate predictor placed midway through the
aggregator emits a per-patch logit $g_j$ for every patch token $j$, which enters the
subsequent global-attention layers as an additive bias $B$ on the attention logits:
\begin{equation}
  A = \operatorname{softmax}\!\Big( \frac{QK^{\top}}{\sqrt{d}} + B \Big),
  \qquad
  B_{ij} =
  \begin{cases}
    \phi(g_j), & i \in \text{CameraTokens},\; j \in \text{PatchTokens} \\[2pt]
    0,         & \text{otherwise,}
  \end{cases}
\end{equation}
where $\phi$ turns the logit into a suppression strength,
\begin{equation}
  \phi(g) =
  \begin{cases}
    x,          & x \le 0 \\
    \lambda x,  & x > 0
  \end{cases}
  \qquad
  x = \operatorname{softplus}(0) - \operatorname{softplus}(g),
  \quad \lambda = 0.1 .
\end{equation}
Only the rows of the camera token carry the bias, so a patch judged dynamic is suppressed
exactly where the pose is read out, while patch-to-patch attention --- and with it the depth
and point predictions --- is left untouched. The construction has two properties we rely on:
$\phi(0) = 0$, so a freshly initialised gate reproduces the pretrained model exactly, and
$\phi \le 0$ up to the small leak $\lambda$, so the gate can only down-weight a patch, never
boost one. The leak keeps a gradient alive for patches that have drifted to the confident-static
side, which would otherwise never come back. Because $B$ sits inside the attention,
$\mathcal{L}_{\text{camera}}$ back-propagates into $g_j$: the pose error is the gate's only
teacher. This is what makes the method applicable to endoscopy, where no ground-truth dynamic
mask exists, and $\sigma(g_j)$ can be read out as a dynamic mask at no extra cost.

\paragraph{Trajectory smoothness.}
The last term is a regulariser on the predicted trajectory alone and uses no ground truth.
Let $\Delta t_i$ be the number of frames that separate $I_i$ and $I_{i+1}$ in the original
video: the frames of a training clip are not sampled at a fixed stride, so this gap varies and
has to be divided out. The predicted camera velocities are then the first differences of the
same quantities the camera loss supervises, per unit time:
\begin{equation}
  v^{T}_i = \frac{\hat{\mathbf{t}}_{i+1} - \hat{\mathbf{t}}_i}{\Delta t_i},
  \qquad
  v^{R}_i = \frac{\hat{\mathbf{q}}_{i+1} - \hat{\mathbf{q}}_i}{\Delta t_i} .
\end{equation}
Penalising these velocities directly would amount to asking the camera not to move, so we
penalise their change instead --- the camera acceleration,
\begin{equation}
  \Delta v^{T}_i = v^{T}_{i+1} - v^{T}_i,
  \qquad
  \Delta v^{R}_i = v^{R}_{i+1} - v^{R}_i ,
\end{equation}
averaged over the $N-2$ second differences a clip of $N$ frames admits:
\begin{equation}
  \mathcal{L}_{\text{smooth}}
  = \frac{1}{N-2} \sum_{i=1}^{N-2}
    \Big( \big\| \Delta v^{T}_i \big\|_{1} + \big\| \Delta v^{R}_i \big\|_{1} \Big) .
\end{equation}
A constant camera motion is therefore free of penalty and only jerky pose sequences are
suppressed. This matters for endoscopy, where the model already under-predicts the amount of
camera motion, so a term penalising the motion itself would push in the wrong direction. As
with the camera loss, every refinement stage is supervised with the same $\gamma^{K-k}$
weighting. Pairs with $\Delta t_i = 0$, which arise when the sampler draws the same frame twice, are
dropped, as are the differences built on them. Quaternions are hemisphere-aligned before
differencing, since $q$ and $-q$ denote the same rotation, and the field of view is left
unsmoothed.

\section*{Evaluation Metrics}

Both depth and pose are only recoverable up to an unknown scale, so every metric below is
computed after an explicit alignment step; without it the numbers would measure the missing
scale rather than the geometry.

\paragraph{Depth.}
We follow the evaluation protocol of AF-SfMLearner~\cite{shao2022afsfmlearner} and
EndoSfM3D~\cite{zhang2026endosfm3d} on SCARED~\cite{allan2021scared} so that the numbers
are directly comparable to the published tables. Each frame is aligned on its own by the ratio
of medians,
\begin{equation}
  s_i = \frac{\operatorname{median}_{\mathcal{V}_i} D_i}{\operatorname{median}_{\mathcal{V}_i} \hat{D}_i},
  \qquad
  \tilde{D}_i = \operatorname{clip}\big( s_i \hat{D}_i,\; d_{\min},\; d_{\max} \big),
\end{equation}
where $\mathcal{V}_i$ collects the pixels of frame $i$ whose ground truth falls inside the valid
range $(d_{\min}, d_{\max}) = (0.01, 150)$\,mm and both medians are taken over that set. Dropping
the frame index and writing $\operatorname{mean}_{\mathcal{V}}$ for the average over the valid
pixels, we report the four standard measures of~\cite{eigen2014depth} --- the absolute
relative error (AbsRel), the squared relative error (SqRel), the root mean squared error (RMSE)
and the root mean squared logarithmic error ($\text{RMSE}_{\log}$):
\begin{align}
  \text{AbsRel}
    &= \operatorname*{mean}_{\mathcal{V}}
       \frac{\big| \tilde{D} - D \big|}{D},
  &\text{SqRel}
    &= \operatorname*{mean}_{\mathcal{V}}
       \frac{\big( \tilde{D} - D \big)^{2}}{D}, \\[4pt]
  \text{RMSE}
    &= \sqrt{ \operatorname*{mean}_{\mathcal{V}}
       \big( \tilde{D} - D \big)^{2} },
  &\text{RMSE}_{\log}
    &= \sqrt{ \operatorname*{mean}_{\mathcal{V}}
       \big( \log \tilde{D} - \log D \big)^{2} } .
\end{align}
The four emphasise different failure modes. AbsRel is a relative error and is therefore
insensitive to how far away the surface is, which makes it the primary quantity for comparing
across sequences; SqRel squares the numerator and so responds much more strongly to outlier
predictions; RMSE is an absolute error in millimetres and reports the physical accuracy of the
reconstruction; and $\text{RMSE}_{\log}$ is computed in log space, which up-weights the relative
error at short range --- exactly the regime an endoscope operates in. Every test frame is scored
on its own and the per-frame values are averaged without weighting, as in the reference protocol.

\paragraph{Pose.}
Camera trajectories are scored by the absolute trajectory error (ATE)~\cite{sturm2012benchmark},
in millimetres. The predicted world-to-camera extrinsics are inverted and their translation part
taken as the camera centre of each frame. A monocular trajectory carries no absolute scale, so
the two trajectories are brought onto a common scale before the error is measured. For a segment
of $n$ frames, let $\hat{\mathbf{p}}_j$ and $\mathbf{p}_j$ be the predicted and ground-truth
camera centres, both expressed in the coordinate frame of that segment's first camera; the scale
factor is the least-squares fit
\begin{equation}
  s = \arg\min_{s} \sum_{j=1}^{n} \big\| s\hat{\mathbf{p}}_j - \mathbf{p}_j \big\|_{2}^{2}
    = \frac{\sum_{j} \langle \mathbf{p}_j, \hat{\mathbf{p}}_j \rangle}
           {\sum_{j} \langle \hat{\mathbf{p}}_j, \hat{\mathbf{p}}_j \rangle} ,
\end{equation}
and the ATE of the segment is the residual of that fit,
\begin{equation}
  \text{ATE} = \frac{1}{n} \sqrt{ \sum_{j=1}^{n}
    \big\| s\hat{\mathbf{p}}_j - \mathbf{p}_j \big\|_{2}^{2} } .
\end{equation}
ATE answers a single question: once the unrecoverable scale is taken out, how far does the
estimated camera path stray from the true one? Because it is reported in millimetres it reads as
a physical distance rather than a relative score, and because it is a translational measure a
rotation error enters only through the displacement it induces --- which, over a short window, is
exactly the displacement a surgeon would see as the view drifting off target.

This is evaluated on sliding windows of $n = 5$ frames and averaged over every window of the
sequence, the protocol of AF-SfMLearner~\cite{shao2022afsfmlearner}, under which the published
endoscopic pose tables are computed. Two details of that protocol are worth stating: only the scale is fitted --- the
rotation is not aligned, so an error in the predicted rotation shows up directly in the position
residual --- and the residual is divided by $n$ rather than $\sqrt{n}$. Both follow the reference
implementation, and reproducing them exactly is what makes the number comparable to the published
tables. Scoring short windows keeps the metric sensitive to local tracking quality and prevents a
single early error from dominating the whole sequence; the price is that it does not measure
long-range drift, since every window is re-anchored.

\section*{Results}

Table~\ref{tab:scared_pose_depth} isolates the contribution of each component. Adding the gate
predictor to the fine-tuned VGGT baseline lowers the mean ATE from $0.076$ to $0.070$, and adding
the temporal attention together with the smoothness regulariser brings it to $0.060$, a $21\%$
reduction over the baseline. The depth columns move in the same direction --- AbsRel from $0.045$
to $0.043$, RMSE from $4.254$ to $4.090$ --- even though neither component is supervised by a
depth signal, which suggests that a camera token that reads out of the static part of the scene
also stabilises the geometry the depth head sees.

Table~\ref{tab:scared_sota} places the model against published methods on SCARED. Our model
reports the lowest error on all four depth metrics; on pose it is second to
Endo-FASt3r~\cite{zeinoddin2025endofast3r} on both sequences. The comparison is not like for like
and the G.T.\ column says why: every self-supervised entry trains without ground-truth depth or
pose, whereas our model and the two feed-forward baselines are fine-tuned with both. The input
regime differs as well --- the published methods predict depth from a single image, while our
model consumes the sequence.

Table~\ref{tab:runtime} reports the inference cost. Both feed-forward models run at roughly
$8.8$ frames per second on one 64-frame sequence, against $0.14$ for
MonST3R~\cite{zhang2025monst3r}, and the gate predictor and temporal blocks are not measurable
against the VGGT baseline.

\begin{table}[htbp]
  \centering
  \setlength{\tabcolsep}{6pt}
  \renewcommand{\arraystretch}{1.3}
  \resizebox{\textwidth}{!}{%
  \begin{tabular}{l cccc @{\hspace{2.5em}} ccc}
    \toprule
    \multirow{2}{*}{Method}
      & \multicolumn{4}{c}{Depth}
      & \multicolumn{3}{c}{Camera Pose} \\
    \cmidrule(lr){2-5} \cmidrule(lr){6-8}
    & AbsRel$\downarrow$ & SqRel$\downarrow$ & RMSE$\downarrow$ & RMSELog$\downarrow$
      & ATE$\downarrow$ (Seq.\ 1) & ATE$\downarrow$ (Seq.\ 2) & Mean$\downarrow$ \\
    \midrule
    VGGT~\cite{wang2025vggt} & 0.045 & 0.331 & 4.254 & 0.068 & 0.090 & 0.063 & 0.076 \\
    + Gate predictor & 0.045 & 0.318 & 4.226 & 0.066 & 0.083 & 0.057 & 0.070 \\
    \makecell[l]{+ Temporal attention\\and Smoothness loss}
      & \textbf{0.043} & \textbf{0.309} & \textbf{4.090} & \textbf{0.065}
      & \textbf{0.072} & \textbf{0.048} & \textbf{0.060} \\
    \bottomrule
  \end{tabular}%
  }
  \caption{Ablation of the proposed components, all fine-tuned on SCARED. The left half
  reports depth errors, the right half ATE.}
  \label{tab:scared_pose_depth}
\end{table}

\begin{table}[htbp]
  \centering
  \setlength{\tabcolsep}{6pt}
  \renewcommand{\arraystretch}{1.2}
  \resizebox{\textwidth}{!}{%
  \begin{tabular}{l l c cccc @{\hspace{2.5em}} cc}
    \toprule
    \multirow{2}{*}{Method} & \multirow{2}{*}{Venue} & \multirow{2}{*}{G.T.}
      & \multicolumn{4}{c}{Depth}
      & \multicolumn{2}{c}{Camera Pose} \\
    \cmidrule(lr){4-7} \cmidrule(lr){8-9}
    & & & AbsRel$\downarrow$ & SqRel$\downarrow$ & RMSE$\downarrow$ & RMSELog$\downarrow$
      & ATE$\downarrow$ (Seq.\ 1) & ATE$\downarrow$ (Seq.\ 2) \\
    \midrule
    Fang et al.~\cite{fang2020goodpractice} & WACV'20    & $\times$     & 0.078 & 0.794 & 6.794 & 0.109 & ---   & ---   \\
    Monodepth2~\cite{godard2019monodepth2} & ICCV'19    & $\times$     & 0.069 & 0.577 & 5.546 & 0.094 & 0.077 & 0.055 \\
    Endo-SfM~\cite{ozyoruk2021endosfm} & MIA'21     & $\times$     & 0.062 & 0.606 & 5.726 & 0.093 & 0.076 & 0.050 \\
    AF-SfMLearner~\cite{shao2022afsfmlearner} & MIA'22     & $\times$     & 0.059 & 0.435 & 4.925 & 0.082 & 0.074 & 0.048 \\
    DARES~\cite{zeinoddin2024dares} & ECCVW'24   & $\times$     & 0.052 & 0.356 & 4.483 & 0.073 & ---   & ---   \\
    EndoDAC~\cite{cui2024endodac} & MICCAI'24  & $\times$     & 0.052 & 0.362 & 4.464 & 0.073 & ---   & ---   \\
    Endo-FASt3r~\cite{zeinoddin2025endofast3r} & MICCAI'25  & $\times$     & 0.051 & 0.354 & 4.480 & ---   & \textbf{0.070} & \textbf{0.044} \\
    EndoSfM3D~\cite{zhang2026endosfm3d} & MICCAIW'25 & $\times$     & 0.050 & 0.389 & 4.749 & 0.070 & 0.079 & 0.053 \\
    ESRT-Full~\cite{chang2026esrt} & MICCAI'26  & $\times$     & 0.053 & 0.454 & 4.791 & 0.076 & ---   & ---   \\
    ESRT-Row~\cite{chang2026esrt} & MICCAI'26  & $\times$     & 0.048 & 0.372 & 4.242 & 0.069 & ---   & ---   \\
    DUSt3R~\cite{wang2024dust3r} & CVPR'24    & $\checkmark$ & 0.059 & 0.516 & 4.947 & 0.080 & ---   & ---   \\
    \midrule
    VGGT~\cite{wang2025vggt} & --- & $\checkmark$ & 0.045 & 0.331 & 4.254 & 0.068 & 0.090 & 0.063 \\
    Ours & --- & $\checkmark$ & \textbf{0.043} & \textbf{0.309} & \textbf{4.090} & \textbf{0.065} & 0.072 & 0.048 \\
    \bottomrule
  \end{tabular}%
  }
  \caption{Comparison against published methods on SCARED. Depth follows the
  AF-SfMLearner~\cite{shao2022afsfmlearner} protocol; pose is the 5-frame snippet ATE, with Seq.~1 = \texttt{dataset5/keyframe4}
  (411 frames) and Seq.~2 = \texttt{dataset3/keyframe4} (834 frames). A dash marks a metric the
  corresponding paper does not report; the G.T.\ column marks whether the method is
  trained with ground-truth depth and pose ($\checkmark$) or self-supervised ($\times$). Published methods take a single image (DUSt3R and ESRT a
  pair); our baseline and our model take the full sequence for depth and a 64-frame context for
  pose.}
  \label{tab:scared_sota}
\end{table}

\begin{table}[htbp]
  \centering
  \setlength{\tabcolsep}{8pt}
  \renewcommand{\arraystretch}{1.2}
  \begin{tabular}{l c r r}
    \toprule
    Method & Inference & Throughput (FPS)$\uparrow$ & Time per frame (ms)$\downarrow$ \\
    \midrule
    MonST3R~\cite{zhang2025monst3r} & \makecell{optical flow +\\global alignment} &   0.14 & 7250.2 \\
    VGGT~\cite{wang2025vggt} & feed-forward &   8.71 &  114.8 \\
    Ours    & feed-forward &   8.81 &  113.5 \\
    \bottomrule
  \end{tabular}
  \caption{Inference cost on one 64-frame sequence (\texttt{dataset2/keyframe3}), measured on a
  single RTX 5090 under identical conditions: image loading and preprocessing are included,
  model construction and any writing of results are not. MonST3R~\cite{zhang2025monst3r} is two
  orders of magnitude
  slower because it precomputes RAFT optical flow and then runs 300 iterations of global
  alignment per sequence, whereas both feed-forward models produce the whole sequence in one
  pass. The gate predictor and the temporal blocks add no measurable cost over the VGGT
  baseline. These are throughput figures, not streaming latency: all three methods consume the
  sequence as a whole, so no frame is available before the sequence finishes.}
  \label{tab:runtime}
\end{table}

\bibliographystyle{IEEEtran}
\bibliography{references}

\end{document}
